\infographicchapter{Optical Coherence Tomography Angiography}{octa.png}
\label{ch:infographic-octa}

\section{Overview}

Optical coherence tomography angiography (OCTA) is a fast, painless, dye-free eye scan that maps the tiny blood vessels at the back of your eye in three dimensions. It builds on standard OCT by taking several rapid pictures of the exact same spot on your retina within fractions of a second. Stationary eye tissues look identical from one snapshot to the next, while red blood cells moving through blood vessels cause tiny flicker-like changes. The scanner's computer isolates that motion to draw a sharp, layered map of active blood flow---no needles or dye required.

OCTA separates the back of the eye into thin, individual layers, allowing eye doctors to view both upper and deeper retinal capillary networks, the normally blood-vessel-free outer retina, and the rich blood layer beneath it (the choriocapillaris). Doctors commonly use it to evaluate diabetic eye disease (checking for blocked capillaries and widened blood-free zones), spot abnormal new blood vessels in wet age-related macular degeneration, assess blocked retinal veins or arteries, evaluate rare macular vessel conditions, and check tiny vessels around the optic nerve in glaucoma.

Compared with traditional dye angiography, OCTA avoids IV needles, nausea, and rare allergic reactions, takes only seconds, and never has leaking dye that blurs deeper layers. Its main limits are that it cannot show active dye leakage or very sluggish blood flow, it focuses mostly on the central retina rather than the far edges, and eye movement, blinking, or severely swollen retinas can create shadows or tracing errors that doctors must carefully review.

\section{Underlying Physical Principles}
\begin{itemize}
\item \textbf{Ultra-Fast Light Scanning:} OCTA uses safe, high-speed near-infrared light (taking 70,000 to 400,000 tiny depth measurements every second) to snap several back-to-back cross-sectional pictures at the exact same spot on your retina within milliseconds.
\item \textbf{Tracking Moving Red Blood Cells:} Still retinal layers reflect light steadily, whereas red blood cells flowing through tiny capillaries change the reflected light pattern from moment to moment. This natural movement creates blood-vessel contrast without any injected dye.
\item \textbf{Smart Motion-Detection Software:} Specialized algorithms---such as split-spectrum amplitude-decorrelation (SSADA)---measure signal changes across repeat cross-sectional slices (B-scans) and split the light spectrum into narrower bands to boost capillary clarity while filtering out gentle whole-eye motion, highlighting vessels where blood flows above a minimum speed.
\item \textbf{Built-In 3D Alignment:} Because the flow map comes from the exact same light scan as structural OCT, every tiny blood vessel lines up automatically with the exact retinal layer it sits in.
\item \textbf{Layer-by-Layer Vessel Maps:} Software slices the 3D scan into clean flat layers (upper capillaries, deep capillaries, outer retina, and underlying choroid), cleans up shadows from surface vessels, and calculates vessel density and central blood-free zone size.
\end{itemize}

\section{Imaging Mechanism}
\begin{enumerate}
\item \textbf{Positioning and Eye-Tracking Lock:} You sit comfortably with your chin and forehead resting on the scanner's supports and focus on a small target light. Built-in infrared eye tracking locks onto your retina to compensate for tiny natural eye shifts.
\item \textbf{Repeated Rapid Sweeps:} Tiny mirrors sweep the gentle near-infrared beam across a small grid (typically 3x3~mm to 12x12~mm) on your retina, capturing 2 to 4 rapid repeat slices at each exact position.
\item \textbf{High-Speed Digital Capture:} Reflected light recombines with an internal reference beam and hits an ultra-fast digital sensor that records hundreds of thousands of measurements per second.
\item \textbf{Extracting Blood Flow and Removing Shadows:} The computer builds standard structural slices, compares repeat scans to map moving red blood cells, separates the retina into anatomical layers, and applies filters to remove shadow projections from upper vessels.
\item \textbf{Displaying Layered Vessel Maps:} The workstation displays crisp top-down (en-face) blood-vessel maps alongside color-coded flow overlays on cross-sectional slices and automated vessel-density numbers.
\end{enumerate}

\section{Key Technical Parameters}
\noindent
\begin{tabularx}{\textwidth}{@{}l X@{}}
\toprule
\textbf{Parameter} & \textbf{Typical Value / Description} \\
\midrule
Physical Principle & Detecting natural red-blood-cell movement across rapid repeat near-infrared light scans at the exact same retinal spot \\
Energy / Signal Source & Safe, low-power Class~1 near-infrared light (SD-OCTA or high-speed swept-source SS-OCTA) \\
Typical Frequency / Energy & 840~nm near-infrared light (SD-OCTA); 1050~nm (SS-OCTA for deeper views) \\
Spatial Resolution & 5--7~micrometers in depth; 10--15~micrometers side-to-side (standard 3x3~mm to 6x6~mm, up to 12x12~mm widefield) \\
Typical Scan Duration & 3--10~seconds per eye capture; 5--15~minutes total visit for both eyes \\
Contrast Agent & None (completely dye-free; tracks natural blood-cell motion) \\
Ionizing Radiation & None (0~mSv; safe non-ionizing near-infrared light) \\
\bottomrule
\end{tabularx}

\section{Common Symptoms That Prompt This Scan}
\begin{itemize}
\item \textbf{Blurry or Distorted Central Vision:} Wavy straight lines or fuzzy reading vision when eye doctors want a dye-free 3D map of abnormal leaking blood vessels in macular degeneration.
\item \textbf{Diabetic Eye Checkups (Often Without Symptoms):} Spotting tiny areas of poor retinal blood flow and fragile new vessels early in people with diabetes without needing an IV dye injection.
\item \textbf{Sudden Vision Loss From a Blocked Retinal Vein:} Measuring how much retinal circulation was interrupted and monitoring recovery over repeat clinic visits.
\item \textbf{Monitoring Eye Injections Safely Over Time:} Tracking whether abnormal retinal vessels are shrinking after anti-VEGF treatments without repeated needle dyes.
\end{itemize}

\section{Clinical Usages}

\subsection{Primary Usages}
\begin{itemize}
\item \textbf{Age-Related Macular Degeneration (Dry and Wet AMD):} Spots abnormal new blood vessels growing under or into the macula in wet AMD---including quiet, early vessel networks before they start leaking---and tracks deeper choriocapillaris flow loss around patches of thinning tissue (geographic atrophy) in advanced dry AMD, all without needing an IV dye injection.
\item \textbf{Diabetic Eye Disease (Diabetic Retinopathy):} Maps areas where tiny capillaries have shut down, measures widening of the central blood-free zone (foveal avascular zone), and highlights abnormal fragile vessels.
\item \textbf{Blocked Retinal Veins and Arteries:} Shows how much circulation has been lost in the macula, where capillaries have dropped out, and where natural bypass vessels have formed.
\item \textbf{Macular Telangiectasia and Chronic Fluid Disorders:} Identifies abnormal dilated capillaries in macular telangiectasia type~2 and spots hidden new blood vessels complicating chronic fluid under the retina (central serous chorioretinopathy).
\item \textbf{Glaucoma Blood-Flow Checks:} Measures the density of tiny capillaries around the optic nerve and macula to help detect early circulation loss alongside standard optic nerve thickness checks.
\end{itemize}

\subsection{Secondary Usages}
\begin{itemize}
\item \textbf{Tracking Eye Injection (Anti-VEGF) Response:} Watches how abnormal blood vessels shrink, reshape, or return during anti-VEGF treatment without repeating dye injections.
\item \textbf{Eye Inflammation (Uveitis):} Detects abnormal new vessels and patchy blood-flow gaps beneath the retina caused by inflammatory eye conditions.
\item \textbf{High Nearsightedness (Myopic Maculopathy):} Quickly tells abnormal leaking-prone vessels apart from harmless tissue cracks or small bleeds in very nearsighted eyes.
\item \textbf{Sickle Cell and Radiation Retinopathy:} Detects early macular capillary drop-out caused by sickle cell disease or prior radiation therapy near the eye, often before noticeable vision changes appear.
\item \textbf{Optic Nerve and Research Uses:} Assists in checking optic nerve blood supply in nerve strokes (ischemic optic neuropathy) and serves as a research tool for studying brain and whole-body blood-vessel health.
\end{itemize}

\section{Advantages}
\begin{itemize}
\item \textbf{Completely Noninvasive and Dye-Free:} Needs no IV needle or fluorescent dye, avoiding nausea, hives, IV leaks, and rare severe allergic reactions---and making it completely safe during pregnancy and kidney disease.
\item \textbf{Clear Layer-by-Layer Views:} Separates upper capillaries, deep capillaries, the outer retina, and the deeper choroid into distinct maps, revealing hidden deep-vessel problems that traditional 2D dye tests blur together.
\item \textbf{No Blurring from Leaking Dye:} Because no dye leaks out into surrounding tissue, delicate abnormal vessel webs and areas of poor circulation stay sharp and easy to see.
\item \textbf{Fast Enough for Every Clinic Visit:} Takes only seconds per eye, allowing doctors to track vessel changes and treatment response regularly.
\item \textbf{Built-In Measurements:} Automatically calculates vessel density and the size of the central blood-free zone while matching them directly to structural OCT slices.
\end{itemize}

\section{Disadvantages}
\begin{itemize}
\item \textbf{Cannot Show Active Dye Leakage:} Because OCTA only tracks moving blood cells, it cannot show active leaking dye, stained vessel walls, or fluid pooling needed in certain inflammatory or leaking eye diseases.
\item \textbf{Can Miss Very Slow Blood Flow:} Blood moving more slowly than the scanner's detection limit---such as inside some tiny microaneurysms, polyp-like vessels, or scarred vessel webs---may not show up.
\item \textbf{Shadow and Movement Glitches:} Upper blood vessels can cast false ``shadow copies'' onto deeper layers (mimicking abnormal vessels), and blinking or tiny eye jumps can create straight-line streaks.
\item \textbf{Layer-Tracing Errors in Swollen Retinas:} Heavy macular swelling, lifted retinal layers, or high nearsightedness can confuse the automated layer borders, requiring careful manual review by the eye doctor.
\item \textbf{Narrow View and Blocked by Cloudiness:} Standard scans (3x3~mm to 6x6~mm) focus on the center of the retina and can miss far-edge circulation loss unless multiple scans are stitched together. Thick cataracts or eye bleeding also weaken the image.
\end{itemize}

\section{Contraindications}

\subsection{Absolute}
\begin{itemize}
\item \textbf{None (Standard Undilated Exam):} Standard undilated OCTA has no medical contraindications. Because no dye or radiation is used, it is safe during pregnancy, kidney failure, and severe fluorescein or iodine allergies.
\end{itemize}

\subsection{Relative}
\begin{itemize}
\item \textbf{Dilating Eye Drops (When Needed):} If dilating drops are needed to widen a small pupil, doctors avoid them in eyes with very narrow drainage angles (risk of acute angle-closure glaucoma) or known allergies to dilating drops.
\item \textbf{Eye Movement, Infections, and Cloudiness:} Severe involuntary eye shaking (nystagmus), trouble focusing steadily, contagious pink eye, or dense cataracts, corneal scars, and heavy eye bleeding can prevent a clear scan.
\end{itemize}

\section{Patient Safety Guidelines}
\begin{itemize}
\item \textbf{Safe Light Exposure:} Clinics use calibrated Class~1 eye-safe OCTA machines (meeting ANSI~Z136.1 and IEC~60825-1 standards) and avoid unnecessary extra repeat scans over the same central spot.
\item \textbf{Clean Equipment Between Patients:} Staff disinfect the chin rest, forehead bar, and handgrips after every patient and reschedule non-urgent scans if contagious pink eye is suspected.
\item \textbf{Safe Dilation Practices:} When dilating drops are used for small pupils or cataracts, staff check your drainage angle and allergy history first and advise you about temporary glare and driving safely.
\item \textbf{Eye Comfort and Rest Breaks:} Because OCTA requires holding your gaze steady for several seconds, staff offer short blink breaks, preservative-free artificial tears for dry eyes, and comfortable chair support.
\item \textbf{Knowing When Dye Testing Is Still Needed:} While OCTA safely avoids IV dye risks in pregnancy, kidney disease, and dye allergies, doctors will still recommend traditional dye angiography (with full safety checks) if seeing active vessel leakage is essential.
\end{itemize}

\section{Prerequisite Tests}
\begin{itemize}
\item \textbf{No Blood Tests or Fasting Needed:} Because OCTA is completely dye-free and noninvasive, no blood work, kidney panels, or pregnancy tests are needed beforehand.
\item \textbf{Baseline Eye Exam:} Your eye doctor checks your vision, inspects the front of your eye, and looks at your retina to confirm the reason for the scan and spot dry eye, cataracts, or bleeding that could affect clarity.
\item \textbf{Standard Structural OCT:} Regular cross-sectional OCT is almost always taken at the same time on the same machine so the doctor can check retinal swelling and make sure automated vessel layers are traced accurately.
\item \textbf{Angle and Focus Check:} If dilating drops are planned, your doctor checks the front drainage angle and allergy history beforehand, and sees how steadily you can focus on a target.
\end{itemize}

\section{Preparing for the Scan}
\begin{itemize}
\item \textbf{Take Out Contact Lenses:} Wear your glasses or bring a clean contact lens case and solution, as contact lenses are removed beforehand to keep the optical view crisp and prevent surface dryness.
\item \textbf{Ask About Dilating Drops:} Many central OCTA scans can be taken without dilating drops, though dilation is still used for small pupils, cataracts, wide-area scans, or a full retinal exam. Check ahead of time, and if drops are planned, bring dark sunglasses and a companion to drive you home.
\item \textbf{Normal Meals and Medicines:} Because no IV dye is used, you can eat normally and take all regular pills and prescribed eye drops (including glaucoma drops). Just skip thick eye ointments right before your visit.
\item \textbf{Lubricate Dry Eyes:} Tell the technician if your eyes feel dry or gritty so preservative-free artificial tears can be added first. Adjust your chair so your neck feels relaxed with your chin and forehead squarely on the supports.
\item \textbf{Steady Gaze Practice:} OCTA needs a few extra seconds of stillness to detect moving blood cells. Take one clean blink right before the scan begins, then lock your eyes gently on the target light without squinting or moving your head.
\end{itemize}

\section{During Scan}
\begin{itemize}
\item \textbf{Holding Your Head Steady:} You will sit upright at the scanner with your chin firmly in the chin rest and your forehead touching the top bar. Try not to lean back or shift during each quick capture.
\item \textbf{Blinking and Breathing:} Breathe calmly and naturally. Take a full blink right before the technician starts each capture, then keep your eyelids comfortably open without squeezing for a few seconds.
\item \textbf{What You Will Notice:} Keep looking steadily at the small target light inside the lens while the painless near-infrared beam sweeps across your eye. There are no needles, no dye, no bright camera flashes, and nothing touches your eye.
\item \textbf{Talking with the Technician:} The technician sits right across from you and guides each step. Let them know immediately if your eyes feel dry, watery, or tired so you can pause, blink, or get a soothing lubricating drop.
\item \textbf{Scan Time:} Each eye capture takes about 3--10~seconds (just a touch longer than standard OCT), and a full scan of both eyes usually finishes in 5--15~minutes.
\end{itemize}

\section{Normal Appearance}
\begin{itemize}
\item \textbf{Rich, Even Capillary Web:} Normal upper and deep retinal maps look like a fine, lacy, even spiderweb of tiny blood vessels with no dark empty patches or balloon-like bulges.
\item \textbf{Smooth Central Blood-Free Zone (FAZ):} Right at the center of the macula sits a small, clean, round-to-oval dark circle (typically 0.2--0.4~square~mm) surrounded by an unbroken ring of tiny capillaries.
\item \textbf{Clear Outer Retina:} The deeper outer retinal layer normally contains no blood vessels, so it appears smoothly dark on a healthy scan.
\item \textbf{Even Deeper Blood Layer and Optic Nerve Ring:} The layer beneath the retina (choriocapillaris) looks like an even, fine ``salt-and-pepper'' texture without dark gaps, and fine capillaries fan out evenly around the optic nerve.
\end{itemize}

\section{Observing for Abnormal Patterns}
\begin{itemize}
\item \textbf{Abnormal New Blood Vessels (Macular Neovascularization):} Tangled, branching, or fan-shaped webs of bright blood flow appearing inside the normally dark outer retina---confirmed on cross-sectional OCT slices---identify abnormal new vessels in wet AMD, high nearsightedness, or chronic fluid disorders.
\item \textbf{Missing Capillaries and Widened Central Zone:} Dark, empty patches where the lacy capillary web has disappeared indicate poor blood supply (ischemia). A widened or ragged central blood-free zone shows macular circulation loss in diabetic eye disease or blocked retinal veins.
\item \textbf{Vessel Bulges and Fragile Growths:} Tiny dot-like bulges (microaneurysms), twisted bypass vessels, and fragile new vessels growing forward into the eye's gel mark advancing diabetic retinopathy.
\item \textbf{Deeper Flow Gaps and Optic Nerve Loss:} Patchy dark gaps in the normally even deeper blood layer point to inflammation or tissue thinning, while missing wedges of capillaries around the optic nerve often match glaucoma vision loss.
\item \textbf{False Look-Alikes (Artifacts):} Upper blood vessels can cast shadow copies onto deeper layers that look like abnormal new vessels; heavy swelling can misplace normal vessels into the wrong layer; and blood, drusen, or very slow flow can create dark spots that mimic missing circulation.
\end{itemize}

\section{Dealing with Uncertainties}
\begin{itemize}
\item \textbf{Upper-Vessel Shadows and Eye Movement:} Shadows from upper vessels can mimic abnormal vessels in deeper layers, while quick eye jumps create horizontal bands or doubled vessels. Doctors turn on shadow-removal software, check cross-sectional flow overlays, and repeat scans with steady focus.
\item \textbf{Layer-Tracing Mix-Ups and Blocked Light:} Heavy swelling, lifted retinal layers, or high nearsightedness can shift automated layer borders, while bleeding, drusen, or cataracts block light and mimic missing vessels. Doctors manually adjust layer lines and compare with structural OCT slices.
\item \textbf{Very Slow Flow and Edge Blind Spots:} Very slow blood flow (in some microaneurysms, polyp-like vessels, or scarred webs) can fall below OCTA's detection limit, and standard 3x3~mm to 6x6~mm scans do not reach the far edges of the retina.
\item \textbf{Confirming with Dye Angiography:} Because OCTA cannot show active dye leakage or staining, doctors may still order traditional fluorescein or indocyanine green dye angiography---alongside a dilated eye exam---when active leaking, eye inflammation, or far-edge circulation loss must be confirmed.
\end{itemize}

\section{Augmenting Methods}
\begin{itemize}
\item \textbf{Paired Structural OCT Slices:} Viewing OCTA right beside standard cross-sectional OCT with color flow overlays confirms retinal layer health, spots fluid pockets, and separates real abnormal vessels from shadows or tracing errors.
\item \textbf{Traditional Dye Angiography (FA and ICGA):} Complements OCTA by showing active fluid leakage, inflamed vessel walls, slow-flow bulges, polyp-like choroidal vessels, and wide-angle far-edge circulation loss.
\item \textbf{Wide-Angle Color Photos and Autofluorescence:} Wide-angle photos document bleeding and peripheral vessel changes outside the OCTA window, while autofluorescence maps the health of the retina's support layer.
\item \textbf{Visual Field and General Health Checks:} Visual field testing matches optic nerve vessel loss to side vision in glaucoma, while blood sugar (HbA1c), blood pressure, carotid artery checks, or inflammation blood tests address underlying causes.
\item \textbf{Widefield Stitching and AI Tools:} Stitching multiple OCTA scans together widens the view, while automated density, central-zone, and shadow-removal software make visit-to-visit tracking more reliable.
\end{itemize}

\section{Comparison with Alternative Modalities}
\begin{itemize}
\item \textbf{Versus Fluorescein Dye Angiography (FA):} OCTA is dye-free, painless, repeatable at every visit, and separates deep capillary layers cleanly without dye blur. FA is still better for showing active leaking dye, vessel-wall staining, slow-flow microaneurysms, and far-edge circulation loss.
\item \textbf{Versus Indocyanine Green Angiography (ICGA):} Swept-source OCTA shows deep abnormal vessel webs without IV dye, while ICGA remains superior for polyp-like choroidal vessels, deep leaky choroidal tissue, and deep inflammatory conditions.
\item \textbf{Versus Standard Structural OCT:} Structural OCT excels at measuring retinal thickness and fluid pockets, whereas OCTA directly maps blood flow, missing capillaries, and quiet abnormal vessels before they leak.
\end{itemize}

\section{Follow-up Steps}
\begin{itemize}
\item \textbf{Heading Home After the Scan:} You can leave right away after dye-free OCTA. If dilating drops were used, wear sunglasses for glare, avoid driving until near vision clears, and watch for warning signs of sudden eye pressure spikes (severe pain, halos, or nausea). Preservative-free artificial tears help if your eyes feel dry from holding still.
\item \textbf{Careful Doctor Review:} Your eye specialist always reviews the top-down vessel maps alongside cross-sectional structural OCT slices and shadow-removal tools to be certain abnormal findings are real.
\item \textbf{Adding Dye Testing Only When Needed:} If your doctor needs to see active leaking dye, inflamed vessels, polyp-like lesions, or far-edge retinal circulation beyond the OCTA view, they will schedule traditional FA or ICGA dye testing.
\item \textbf{Easy Repeat Monitoring:} During active eye injection (anti-VEGF) care, OCTA is often repeated every 4--12~weeks to watch abnormal vessels shrink or catch early diabetic circulation changes---all without repeated IV dye injections.
\end{itemize}

\section{Key References}
\begin{itemize}
\item Spaide RF, Fujimoto JG, Waheed NK, Sadda SR, Staurenghi G. Optical coherence tomography angiography. \textit{Progress in Retinal and Eye Research}. 2018;64:1--55.
\item American Academy of Ophthalmology (AAO). \textit{Preferred Practice Pattern Guidelines: Age-Related Macular Degeneration and Diabetic Retinopathy}. San Francisco, CA: AAO.
\item de Carlo TE, Romano A, Waheed NK, Duker JS. A review of optical coherence tomography angiography (OCTA). \textit{International Journal of Retina and Vitreous}. 2015;1:5.
\item Kashani AH, Chen CL, Gahm JK, et al. Optical coherence tomography angiography: a comprehensive review of current methods and clinical applications. \textit{Progress in Retinal and Eye Research}. 2017;60:66--100.
\item Le PH, Kaur K, Patel BC. \textit{Optical Coherence Tomography Angiography}. StatPearls, National Library of Medicine (NIH). Available at: \url{https://www.ncbi.nlm.nih.gov/books/NBK563235/}.
\end{itemize}

\section{Further Reading}
\begin{itemize}
  \item Wikipedia contributors. \textit{Optical coherence tomography angiography}. Wikipedia, The Free Encyclopedia. Available at: \url{https://en.wikipedia.org/wiki/Optical_coherence_tomography_angiography}.
\end{itemize}
